% Ad Atticum 11.23 --- headnote
% ad-atticum-11-23, 9 July 47 BC, Brundisium

\textit{Cicero to Atticus, written from Brundisium on
the seventh day before the Ides of Quintilis 47 BC ---
9 July (the manuscript dateline: \textit{Scr.\ Brundisi
vii Id.\ Quint.\ a.\ 707 (47)}). Camillus has sent
word that he and Atticus have already conferred on
the business Cicero entrusted to them, and Cicero
was hoping for confirmation from Atticus himself ---
though he concedes that there is no real altering the
situation. The greater anxiety is for Atticus's health,
which his last letter had mentioned. A traveller named
Agusius has arrived from Rhodes with intelligence:
young Quintus set out for Caesar on 29 May, and
Philotimus reached Rhodes the day before with letters
for Cicero. Quintus the brother has sent the elder
Cicero what reads as triumphant congratulations ---
on what, Cicero will not say, but the bitterness in
the sentence ``given how grave my error has been, I
cannot in even the most distant thought reach anything
I could find bearable'' tells the whole story.}

\textit{The third paragraph is the most painful piece
of Book 11: Cicero opens the question of Tullia's
divorce from Dolabella. He blames himself and Atticus
for not having acted earlier; ``nothing among our
miseries would have been better than a divorce.''
He runs through the grounds that could have been used
--- Dolabella's debt legislation (\textit{tabulae
novae}), the nocturnal break-ins, the affair with
Metella, the general inventory of his crimes --- and
notes that they would not even have lost the dowry
had they moved in time. Now Dolabella himself seems
to be ``giving us notice'' (the rumour of a statue
to Clodius is the trigger), and Cicero declares his
decision, with Atticus's agreement, that the
\textit{nuntium remitti} --- a formal notice of divorce
--- should be sent. The only open question is timing,
relative to the third installment of the dowry the
husband may sue to recover. The letter closes with
Cicero's near-desperate offer to travel by night,
if necessary, simply to see Atticus.}
